% Einheit 1 von 3 – Unabhängigkeit prüfen: die Produktregel P(A ∩ B) = P(A) · P(B) an Anteilen oder absoluten Häufigkeiten prüfen (Tafel oder Anzahlen erst beschaffen (bank/unabhaengigkeit/e1.jsonl, zusammenbau v0.1)
%% TODO Zeitmarke Einheit 1: Marken-Zeile nicht lesbar
%% TODO Zweigzeile Teil 1: was hier gelernt wird (Satz in Schülersprache aus dem Einheitstitel) – Einheit 1
\einheitenkopf[e1]{Einheit 1 von 3 · Unabhängigkeit prüfen: die Produktregel P(A ∩ B) = P(A) · P(B) an Anteilen oder absoluten Häufigkeiten prüfen (Tafel oder Anzahlen erst beschaffen}
\zweigzeile{Abitur GK}
\setcounter{aufgabe}{8}

%% TODO Titel: Ich-kann-Satz fehlt in der Bank (Platzhalter: Kettenname) – Nr. 9
%% TODO Anweisung: ein Satz über der ersten Teilaufgabe fehlt in der Bank – Nr. 9
\begin{aufgabe}{„Prüfen oder nutzen?“}
\begin{teile}
\teil Aufgabe: „Untersuche, ob die Ereignisse A und B stochastisch unabhängig sind.“ Was ist zu tun? \\ \kreuz{prüfen – die Produktregel anwenden} \\ \kreuz{nutzen – Unabhängigkeit ist vorausgesetzt, eine Gleichung aufstellen}
\end{teile}
\end{aufgabe}

%% TODO Titel: Ich-kann-Satz fehlt in der Bank (Platzhalter: Kettenname) – Nr. 10
%% TODO Anweisung: ein Satz über der ersten Teilaufgabe fehlt in der Bank – Nr. 10
\begin{aufgabe}{Unabhängigkeit prüfen}
%% TODO \rechenplatz{n}: Schrittzahl des Grundfalls fehlt in der Bank (nur bei fester Schreibform, 2.3 b)
\begin{teile}
\teil Eine Vierfeldertafel mit allen Feldern und Rändern liegt vor. Wo stehen die drei Zahlen der Produktregel? \\ \kreuz{alle drei stehen in der Tafel: zwei Ränder und ein Feld} \\ \kreuz{nur die Ränder stehen da, der Schnitt muss berechnet werden} \\ \kreuz{keine steht da, alle drei sind zu beschaffen}
\teil Die Tafel zeigt 100 Schüler nach den Ereignissen M (Mädchen) und F (kommt mit dem Fahrrad). Prüfe mit der Produktregel, ob M und F stochastisch unabhängig sind.

\vierfeldertafel{M}{F}{20,30,50,20,30,50,40,60,100}
\teil In einem Sportverein sind 60 Mitglieder, 36 davon Erwachsene. 20 Mitglieder spielen Volleyball, darunter 8 Jugendliche. Bestimme die fehlenden Anzahlen und prüfe, ob E (Erwachsener) und V (spielt Volleyball) stochastisch unabhängig sind.
\teil In einer Region traten 12\,400 Personen zur Fahrprüfung an, davon waren 3\,100 mindestens 30 Jahre alt. 9\,920 Prüflinge bestanden, darunter 7\,750 jüngere. Prüfe, ob der Anteil der Bestandenen unter den Älteren mit dem Gesamtanteil übereinstimmt, ob Ä (mindestens 30 Jahre) und S (bestanden) stochastisch unabhängig sind, und deute das Ergebnis \hfill (Abitur 2019 GK).
\teil Ein Paketdienst befördert 600 Pakete, 150 davon sind schwer. 90 der schweren Pakete gehen ins Ausland, insgesamt gehen 240 Pakete ins Ausland. Untersuche, ob der Anteil der Auslandspakete unter den schweren Paketen ebenso groß ist wie unter den nicht schweren \hfill (Abitur 2022 GK).
\teil Ein Glücksrad mit vier gleich großen Feldern 1, 2, 3, 4 wird zweimal gedreht. Ereignis A: die Summe der beiden Zahlen ist 5. Ereignis B: beim ersten Drehen erscheint die 4. Prüfe ohne Rechner, ob A und B stochastisch unabhängig sind \hfill (Abitur 2018 LK).
\teil Von den Kunden eines Versandhauses besitzen 70\,\% eine Kundenkarte (Ereignis K). Von den Karteninhabern bestellen 40\,\% online; 15\,\% aller Kunden sind Kunden ohne Karte, die online bestellen. O bezeichnet das Ereignis „bestellt online“. Begründe mithilfe der Ungleichung $0{,}4 \ne 0{,}7 \cdot 0{,}4 + 0{,}15$, dass K und O stochastisch abhängig sind \hfill (Abitur 2022 LK).
\end{teile}
\end{aufgabe}

%% TODO Titel: Ich-kann-Satz fehlt in der Bank (Platzhalter: Kettenname) – Nr. 11
%% TODO Anweisung: ein Satz über der ersten Teilaufgabe fehlt in der Bank – Nr. 11
\begin{aufgabe}{Unabhängigkeit prüfen – Fehler finden, Begründen, Anwendung, Darstellungswechsel}
\begin{teile}
\teil Die Tafel zeigt 100 Schüler nach M (Mädchen) und F (kommt mit dem Fahrrad). Lars sagt: „Es gibt 20 Mädchen mit Fahrrad, der Schnitt ist also nicht leer – M und F sind abhängig.“ Finde den Fehler.

\vierfeldertafel{M}{F}{20,30,50,20,30,50,40,60,100}
\teil Begründe, warum die Prüfung $P(A \cap B) = P(A) \cdot P(B)$ und die Prüfung $P_A(B) = P(B)$ zum selben Ergebnis führen.
\teil In einer Stadt mit 40\,000 Haushalten beziehen 12\,000 Ökostrom. 8\,000 Haushalte besitzen ein Elektroauto, darunter 3\,600 Ökostromkunden. Der Energieversorger will Elektroautobesitzer gezielt für Ökostrom werben. Prüfe, ob „Ökostrom“ und „Elektroauto“ stochastisch unabhängig sind, und sage, ob die gezielte Werbung Sinn hat.
\teil Die Tafel zeigt 100 Schüler nach M (Mädchen) und F (kommt mit dem Fahrrad). Übertrage sie in ein Baumdiagramm mit M in der ersten Stufe und lies an den Ästen der zweiten Stufe ab, ob M und F unabhängig sind.

\vierfeldertafel{M}{F}{20,30,50,20,30,50,40,60,100}
\end{teile}
\end{aufgabe}
